\section{Faithful point truncations of holomorphic germ ideals}

The following local algebra bridge uses the actual ring of holomorphic germs.
It is a supplementary form of Krull separation, distinct from the earlier
torus-coefficient criterion. It identifies membership in an ideal from all
of its point truncations; it does not assert a valuative membership criterion.

\begin{theorem}[Membership detected by every point truncation]
\label{trace:actual-holomorphic-germ-point-truncation-closure}
\uses{lag:germ-noetherian}
\source{demailly-cadg}
\lean{Hironaka.HolomorphicGermPointTruncationClosure.mem_iff_mem_all_point_truncations}
\leanok
Let $\mathcal O_n$ be the ring of holomorphic germs at the origin of
$\mathbf C^n$, for any $n\in\mathbf N$, and let $\mathfrak m$ be its actual
maximal ideal. For every ideal $I\subseteq\mathcal O_n$ and every germ
$f\in\mathcal O_n$,
\[
 f\in I\quad\Longleftrightarrow\quad
 \forall N\in\mathbf N,\quad f\in I+\mathfrak m^N.
\]
The assertion includes $I=0$, $I=\mathcal O_n$, $N=0$, and $n=0$.
No properness, finite presentation, nonzero germ, or ideal-membership
certificate is a hypothesis.
\end{theorem}
\begin{proof}
Write $R=\mathcal O_n$. The genuine holomorphic germ ring is local:
a germ is a unit precisely when its value at the origin is nonzero, and
its nonunits form $\mathfrak m$. The already proved Weierstrass induction
gives Noetherianity of $R$, including the zero-dimensional field case.
The quotient $M=R/I$ is a finite $R$-module, generated by the class of $1$.
Since $\mathfrak m$ is a proper ideal of the Noetherian local ring $R$,
the module form of Krull's intersection theorem gives
\[
 \bigcap_{N\geq0}\mathfrak m^N M=0.
\]
Suppose $f\in I+\mathfrak m^N$ for every $N$. At a fixed $N$, write
$f=a_N+b_N$ with $a_N\in I$ and $b_N\in\mathfrak m^N$. Its class in $M$
equals the class of $b_N$, which is $b_N$ times the class of $1$ and hence
belongs to $\mathfrak m^N M$. This holds for every $N$, so the class of
$f$ is zero by the displayed intersection identity. The kernel of
$R\to R/I$ is exactly $I$, and thus $f\in I$.
Conversely, inclusion of $I$ into each sum $I+\mathfrak m^N$ immediately
gives all truncation memberships. This argument also covers the zero
quotient when $I=R$ and the maximal ideal $0$ when $n=0$.
\end{proof}
